\documentclass[11pt]{article}
\usepackage[margin=1in]{geometry}
\usepackage{amsmath,amssymb,amsthm,mathtools,hyperref}
\newtheorem{theorem}{Theorem}
\newcommand{\per}{\operatorname{per}}
\newcommand{\Sym}{\operatorname{Sym}}
\newcommand{\Fock}{\mathcal F}
\title{Marcus's block-permanent inequality for three blocks}
\author{}
\date{30 October 2026\\[3pt]\small Internally audited research draft; external review pending}
\begin{document}
\maketitle
\begin{abstract}
We prove the three-block case of Marcus's block-permanent inequality for every block size and every complex Hermitian positive semidefinite matrix. A partially transposed permutation operator has an explicit positive decomposition except for its final contraction layer. A single trace inequality absorbs that layer's negative part. Two trace bounds handle the quadratic endpoint, and normal ordering removes the ambient-rank restriction. The conjecture for four or more blocks is not settled here. No equality characterization is asserted.
\end{abstract}
\begin{theorem}\label{main}
Let \(A\) be a complex Hermitian positive semidefinite matrix of order \(3k\), partitioned into nine \(k\times k\) blocks \(A_{ij}\). Set \(G_{ij}=\per(A_{ij})\). Then \(\per(A)\ge\per(G)\), for every \(k\ge1\) and every rank.
\end{theorem}
Marcus's general conjecture allows an arbitrary number \(m\) of blocks \cite{wanless,zhang}. The theorem concerns \(m=3\), with no restriction on block size or rank. We use the complex Bargmann--Fock inner product, conjugate-linear in its first argument, with orthogonal monomials and \(\|z^\alpha\|_{\Fock}^2=\alpha!\).
\begin{theorem}\label{poly}
For arbitrary homogeneous degree-\(k\) polynomials \(P,Q,R\) in \(d\) complex variables,
\[
 \|PQR\|_{\Fock}^2\ge
 \per[\langle P_i,P_j\rangle_{\Fock}]_{i,j=1}^3,
 \qquad(P_1,P_2,P_3)=(P,Q,R),
\]
provided either \(k\ge3\) and \(d\le8k^2/9\), or \(k=2\) and \(d\le4\).
\end{theorem}

\section{A weighted permutation operator}
Put \(V=\mathbb C^d\) and partition \(3k\) tensor positions into three groups of \(k\). For a permutation \(\pi\), let \(M_{ij}(\pi)\) count positions sent from group \(j\) to group \(i\). Its row and column sums are \(k\). Write
\[
 H=S_k\wr S_3,\quad W=\sum_{\pi\notin H}P_\pi,\quad
 h(M)=\#\{(i,j):M_{ij}=k\},
\]
\[
 r_i=2-\sum_j(\mathbf1_{M_{ij}=k}+\mathbf1_{M_{ji}=k}),\qquad
 b_i=\frac{3r_i}{6-2h(M)}.
\]
Set all \(b_i=0\) when \(M\) is \(k\) times a permutation matrix. Entries equal to \(k\) occupy distinct rows and columns, and two force a third. Hence \(h\in\{0,1,3\}\), the denominator vanishes exactly on \(H\), and
\[
                         \sum_i b_i=3\mathbf1_{\pi\notin H}.       \tag{1}
\]
Let \(K_i=\sum_\pi b_i(M(\pi))P_\pi\), and transpose the entire \(i\)-th group to obtain \(\widehat Q_i=K_i^{T_i}\). The coefficients are invariant under internal group permutations. It suffices to prove \(Q=\widehat Q_0\ge0\) on
\[
        \mathcal E=\overline{\Sym^kV}\otimes\Sym^kV\otimes\Sym^kV. \tag{2}
\]
The other \(\widehat Q_i\) follow by relabeling. On a product of three symmetric tensors, the expectation of \(K_i\) is the expectation of \(\widehat Q_i\) after conjugating the \(i\)-th factor. Equation (1) then proves nonnegativity of the expectation of \(W\).

\section{The exact contraction factors}
Let \(Q_\ell\) retain terms with \(M_{00}=k-\ell\), so \(Q=\sum_{\ell=0}^kQ_\ell\). The distinguished row and column are pure when \(\ell=0\), so \(Q_0=0\).
Let \(C_{\ell,p}\), for \(0\le p\le\ell\), contract \(p\) distinguished dual slots with the first ordinary group and \(\ell-p\) with the second. These are literal fixed-slot tensor traces, without differentiation factors. The output degrees are
\((k-\ell,k-p,k-\ell+p)\). Set
\[
 n=2k-\ell,\quad
 E_{\ell,p}=\binom{k}{p}\binom{k}{\ell-p}C_{\ell,p},\quad
 a_\ell=\frac{(k!)^2}{(k-\ell)!}.
\]
Embed all ordinary outputs in \(V^{\otimes n}\). Counting the remaining ordinary permutations gives
\[
                 Q_\ell=a_\ell\sum_{p,q} E_{\ell,p}^*T_{pq}E_{\ell,q}, \tag{3}
\]
with the remaining dual identity understood. Here \(T_{pq}\) is the ordinary permutation sum with the prescribed transport coefficient. The common factor follows by choosing the crossing distinguished slots at both ends, the two crossing bijections, and the uncrossed distinguished permutation:
\[
       \binom{k}{\ell}^{\!2}(\ell!)^2(k-\ell)!
                         =\frac{(k!)^2}{(k-\ell)!}.
\]
The ordinary slot choices give exactly the binomial factors in \(E_{\ell,p}\) and \(E_{\ell,q}\). Each remaining ordinary permutation is counted once.

\section{Every nonfinal layer is positive}
For \(1\le\ell<k\), the distinguished row and column cannot contain an entry \(k\). Thus \(r_0=2\), and at most one ordinary cell can equal \(k\):
\[
 b_0=1+\tfrac12\mathbf1_{\{\text{an ordinary cell equals }k\}}.    \tag{4}
\]
Let \(\Pi_n\) fully symmetrize the ordinary outputs, and let
\(L_\ell=\Pi_n\sum_pE_{\ell,p}\).
For the two pure contractions, reorder outputs so the shortened ordinary group comes first. Their sum, denoted \(B_\ell\), has codomain
\[
 \overline{\Sym^{k-\ell}V}\otimes\Sym^{k-\ell}V\otimes\Sym^kV.
\]
An ordinary entry can equal \(k\) only for boundary types \(p,q\in\{0,\ell\}\). It does so exactly when the uncontracted size-\(k\) input subset is carried to its output counterpart. There are \(k!(k-\ell)!\) such permutations. On internally symmetric outputs they all act as the same reordered identification. Both boundary \(E\)'s have coefficient \(\binom{k}{\ell}\).
The constant term in (4) gives \(n!\Pi_n\). Consequently (3) gives the exact sum of squares
\[
 Q_\ell=\frac{(k!)^2(2k-\ell)!}{(k-\ell)!}L_\ell^*L_\ell
       +\frac{(k!)^3}{2}\binom{k}{\ell}^{\!2}B_\ell^*B_\ell,
       \qquad1\le\ell<k.                                    \tag{5}
\]
This is positive in every ambient dimension.

\section{The final layer is compensated by one trace}
At \(\ell=k\), the ordinary permutation coefficient depends only on whether the types \(p,q\) are boundary types \(0,k\):
\[
\begin{array}{c|ccc}
 (p,q)&\text{both interior}&\text{one boundary}&\text{both boundary}\\ \hline
 b_0&1&3/4&0 .
\end{array}
\]
These cases have respectively zero, one, and three pure cells. Put
\[
 I=\Pi_k\sum_{p=1}^{k-1}\binom{k}{p}^{\!2}C_{k,p},
 \qquad B=B_k=C_{k,0}+C_{k,k}.
\]
Both maps have codomain \(\Sym^kV\). The remaining permutation sum is \(k!\Pi_k\), so
\[
 Q_k=(k!)^3\left[
  (I+\tfrac34B)^*(I+\tfrac34B)-\tfrac9{16}B^*B\right].          \tag{6}
\]
This layer alone need not be positive.

Let \(\tau:\bar V\otimes V\otimes\Sym^kV\to\Sym^kV\) trace the first two factors. Completing the remaining pure contraction gives
\[
            \tau B_{k-1}=B,\qquad \tau\tau^*=d\,\mathrm{Id}.    \tag{7}
\]
Thus \(\|Bx\|^2\le d\|B_{k-1}x\|^2\). Keeping only the second square in (5) at \(\ell=k-1\) and adding (6) gives
\[
 \langle x,Qx\rangle\ge(k!)^3\left[
 \|(I+\tfrac34B)x\|^2+
 \left(\frac{k^2}{2d}-\frac9{16}\right)\|Bx\|^2\right].         \tag{8}
\]
Therefore \(Q\ge0\) for \(d\le8k^2/9\). For every \(k\ge3\), this includes \(d\le2k\).

\section{The quadratic endpoint}
For \(k=2\), put \(C=B_1\), \(C_s=\Pi_3C\), \(C_h=(1-\Pi_3)C\). Equations (5)--(6) read
\[
 Q=112C_s^*C_s+16C_h^*C_h+8I^*I+6(I^*B+B^*I).
\]
Let \(U=\tau C_s\), \(Z=\tau C_h\). Tracing the three ordinary positions in \(\Pi_3C\) gives
\[
 U=(B+I)/3,\quad Z=(2B-I)/3,\qquad B=U+Z,\quad I=2U-Z.
\]
Direct symmetrization of the trace adjoint yields
\[
 \tau\Pi_3\tau^*=\frac{d+2}{3}\mathrm{Id},\qquad
 \tau(1-\Pi_3)\tau^*=\frac{2(d-1)}{3}\mathrm{Id}.
\]
At \(d=4\), both constants are \(2\), and for every input \(x\),
\[
\begin{split}
 \langle x,Qx\rangle
 &=112\|C_sx\|^2+16\|C_hx\|^2+
   56\|Ux\|^2-20\operatorname{Re}\langle Ux,Zx\rangle-4\|Zx\|^2\\
 &\ge112\|Ux\|^2-20\operatorname{Re}\langle Ux,Zx\rangle+4\|Zx\|^2\\
 &=87\|Ux\|^2+4\|Zx-\tfrac52Ux\|^2\ge0 .
\end{split}
\]
Smaller dimensions embed isometrically into dimension four.

\section{Fock transfer and arbitrary rank}
For a symmetric tensor \(a\), use the polynomial
\(P_a(z)=\sum a_{i_1\ldots i_k}z_{i_1}\cdots z_{i_k}\).
Then \(\langle P_a,P_b\rangle_{\Fock}=k!\langle a,b\rangle\).
Expansion of the two permutation sums gives
\[
 \langle a\otimes b\otimes c,W(a\otimes b\otimes c)\rangle
 =\|P_aP_bP_c\|_{\Fock}^2
       -\per[\langle P_i,P_j\rangle_{\Fock}].
\]
The preceding positivity proves Theorem~\ref{poly}.

For completeness, the rank reduction follows from normal ordering:
\[
 \|FT\|_{\Fock}^2=
 \sum_\alpha\frac1{\alpha!}
       \|(\partial^\alpha F)^*(D)T\|_{\Fock}^2.                 \tag{9}
\]
The star conjugates coefficients. This identity follows from the multivariate Leibniz rule for \(F^*(D)M_F\) and the adjointness of multiplication and differentiation. For homogeneous \(F\), the top derivative terms give
\(\|FT\|_{\Fock}^2\ge\|F\|_{\Fock}^2\|T\|_{\Fock}^2\).
For equal positive degrees, the zeroth and top terms are distinct and imply
\[
 \|PQ\|_{\Fock}^2\ge
 \|P\|_{\Fock}^2\|Q\|_{\Fock}^2+|\langle P,Q\rangle_{\Fock}|^2. \tag{10}
\]

Factor \(A\) as a Gram matrix. The three block polynomials \(P,Q,R\), each a product of \(k\) linear forms, satisfy
\(\per A=\|PQR\|_{\Fock}^2\) and \(G_{ij}=\langle P_i,P_j\rangle_{\Fock}\).
Let \(U\) span the vectors in \(P,Q\), so \(\dim U\le2k\).
In orthogonal coordinates \(x\in U,y\in U^\perp\), write
\[
 R(x,y)=\sum_\beta y^\beta R_\beta(x),\qquad
 \rho=\sum_{\beta\ne0}\beta!\|R_\beta\|_{\Fock}^2.
\]
Orthogonality and (9) give
\(\|PQR\|_{\Fock}^2\ge\|PQR_0\|_{\Fock}^2+\rho\|PQ\|_{\Fock}^2\).
Replacing \(R\) by \(R_0\) changes only the final diagonal entry of its Gram matrix, so its permanent decreases by
\(\rho\,\per\operatorname{Gram}(P,Q)\).
The original defect is at least the projected defect plus the nonnegative quantity
\[
            \rho[\|PQ\|_{\Fock}^2-\per\operatorname{Gram}(P,Q)].
\]
Theorem~\ref{poly} applies to the projected problem, and (10) supplies the last inequality. Projection preserves the product-of-linear-forms structure, including zero factors. This proves Theorem~\ref{main} for \(k\ge2\), at every rank. For \(k=1\), \(G=A\).

\section*{Verification and scope}
The contraction counting and compensation were derived independently by two agents and audited by a third. Exact standalone checks verify the transport coefficients, subset-kernel formula and literal trace relation. A separate implementation with Gaussian rational polynomial coefficients checks the complete cyclic layer identity against the Fock permanent defect on 24 complex test triples. These checks corroborate normalization; the all-\(k\) conclusion follows analytically from (3)--(10).

This draft has not received external mathematical review or Lean formalization. A bounded primary-source search found the general conjecture still listed as open, with the two-block theorem and restricted three-block cases recorded \cite{wanless,lavi}. This is preliminary priority review. The note makes no assertion about four or more blocks.

\begin{thebibliography}{9}
\bibitem{wanless} I. M. Wanless, \emph{Lieb's permanental dominance conjecture}, 2022, Section 2, Conjecture 3. \url{https://arxiv.org/abs/2202.01867}.
\bibitem{zhang} F. Zhang, \emph{An update on a few permanent conjectures}, Special Matrices 4 (2016), 305--316. \url{https://doi.org/10.1515/spma-2016-0030}.
\bibitem{lavi} Y. Lavi, \emph{Minc's list}, Conjecture 9, updated 29 October 2026; accessed 30 October 2026. \url{https://yair-lavi.com/mincs-list/}.
\end{thebibliography}
\end{document}